% CHƯƠNG 6: sáu thành phần cố định cho mỗi nội dung; hình/bảng chèn tự động; ô TODO do nhóm điền.
\clearpage
\section[Kiểm chứng dự báo và đánh giá dự báo]{KIỂM CHỨNG DỰ BÁO VÀ ĐÁNH GIÁ DỰ BÁO}

\TODO{Đoạn dẫn: các bộ dữ liệu áp dụng (Nile, AirPassengers, DAX log giá); cách đọc một nội dung; câu nối sang Chương 7.}

\subsection{Hàm dự báo}

\noidung{Ảnh hưởng vĩnh viễn của cú sốc}
\tp{Nội dung lý thuyết cần kiểm chứng} \TODO{một câu mệnh đề cần kiểm chứng; dẫn mục 2.3.1(a)}
\tp{Cách triển khai} Dữ liệu thực. \textit{Bộ áp dụng:} Nile, DAX (log giá). \textit{Không áp dụng cho:} AirPassengers \TODO{lý do ngắn}. \TODO{quy trình cụ thể, tham số riêng của mục này}
\tp{Mã lệnh Python} \TODO{đoạn mã lõi 10--20 dòng (chỉ ký tự ASCII và chữ Việt, không dùng ký hiệu Hy Lạp trong mã); mã đầy đủ ở notebook}
\tp{Kết quả tính toán}
\input{\tabdir/tab_b12_06_psi_weights.tex}
\figx{fig_b12_06_psi_weights}{Ảnh hưởng của một cú sốc đơn vị lên $Y_{t+k}$ ($\Psi_k$)}
\TODO{chọn bảng/hình giữ lại, bỏ bớt phần phụ; viết câu dẫn}
\tp{Phân tích và diễn giải}
\bodl{(i) Nile.} \TODO{diễn giải gắn với 2.3.1(a)}

\bodl{(ii) DAX (log giá).} \TODO{diễn giải gắn với 2.3.1(a)}

\tp{Nhận xét} \TODO{xếp loại (phù hợp / phù hợp một phần / không phù hợp); nếu sai lệch thì nêu nguyên nhân}

\noidung{Làm trơn mũ và mô hình mức cục bộ}
\tp{Nội dung lý thuyết cần kiểm chứng} \TODO{một câu mệnh đề cần kiểm chứng; dẫn mục 2.3.1(b)}
\tp{Cách triển khai} Dữ liệu thực. \textit{Bộ áp dụng:} Nile. \textit{Không áp dụng cho:} AirPassengers, DAX (log giá) \TODO{lý do ngắn}. \TODO{quy trình cụ thể, tham số riêng của mục này}
\tp{Mã lệnh Python} \TODO{đoạn mã lõi 10--20 dòng (chỉ ký tự ASCII và chữ Việt, không dùng ký hiệu Hy Lạp trong mã); mã đầy đủ ở notebook}
\tp{Kết quả tính toán}
\input{\tabdir/tab_b12_07_local_level.tex}
\input{\tabdir/tab_b12_07_ses_forecasts.tex}
\TODO{chọn bảng/hình giữ lại, bỏ bớt phần phụ; viết câu dẫn}
\tp{Phân tích và diễn giải}
\bodl{(i) Nile.} \TODO{diễn giải gắn với 2.3.1(b)}

\tp{Nhận xét} \TODO{xếp loại (phù hợp / phù hợp một phần / không phù hợp); nếu sai lệch thì nêu nguyên nhân}

\subsection{Khoảng dự báo}

\noidung{Dự báo tối ưu và khoảng dự báo}
\tp{Nội dung lý thuyết cần kiểm chứng} \TODO{một câu mệnh đề cần kiểm chứng; dẫn mục 2.3.2(a)}
\tp{Cách triển khai} Cả hai. \textit{Bộ áp dụng:} Nile, AirPassengers, DAX (log giá). \TODO{quy trình cụ thể, tham số riêng của mục này}
\tp{Mã lệnh Python} \TODO{đoạn mã lõi 10--20 dòng (chỉ ký tự ASCII và chữ Việt, không dùng ký hiệu Hy Lạp trong mã); mã đầy đủ ở notebook}
\tp{Kết quả tính toán}
\input{\tabdir/tab_b12_08_coverage.tex}
\input{\tabdir/tab_b12_08_forecast_variance.tex}
\figx{fig_b12_08_ni_forecast}{NI: dự báo ARIMA(0,1,1) và khoảng dự báo 95\% cho 20 năm tiếp theo}
\figx{fig_b12_08_se_growth}{Sai số chuẩn của dự báo $h$ bước ($\sigma=1$): $d=0$ so với $d=1$}
\TODO{chọn bảng/hình giữ lại, bỏ bớt phần phụ; viết câu dẫn}
\tp{Phân tích và diễn giải}
\bodl{(i) Nile.} \TODO{diễn giải gắn với 2.3.2(a)}

\bodl{(ii) AirPassengers.} \TODO{diễn giải gắn với 2.3.2(a)}

\bodl{(iii) DAX (log giá).} \TODO{diễn giải gắn với 2.3.2(a)}

\tp{Nhận xét} \TODO{xếp loại (phù hợp / phù hợp một phần / không phù hợp); nếu sai lệch thì nêu nguyên nhân}

\noidung{Dự báo trên thang logarit}
\tp{Nội dung lý thuyết cần kiểm chứng} \TODO{một câu mệnh đề cần kiểm chứng; dẫn mục 2.3.2(b)}
\tp{Cách triển khai} Dữ liệu thực. \textit{Bộ áp dụng:} AirPassengers, DAX (log giá). \textit{Không áp dụng cho:} Nile \TODO{lý do ngắn}. \TODO{quy trình cụ thể, tham số riêng của mục này}
\tp{Mã lệnh Python} \TODO{đoạn mã lõi 10--20 dòng (chỉ ký tự ASCII và chữ Việt, không dùng ký hiệu Hy Lạp trong mã); mã đầy đủ ở notebook}
\tp{Kết quả tính toán}
\input{\tabdir/tab_b12_09_ap_backtransform.tex}
\input{\tabdir/tab_b12_09_dp_forecast.tex}
\figx{fig_b12_09_ap_forecast}{AP: dự báo airline 1959--1960 trên thang gốc (trung vị, khoảng 95\%) và giá trị thực}
\TODO{chọn bảng/hình giữ lại, bỏ bớt phần phụ; viết câu dẫn}
\tp{Phân tích và diễn giải}
\bodl{(i) AirPassengers.} \TODO{diễn giải gắn với 2.3.2(b)}

\bodl{(ii) DAX (log giá).} \TODO{diễn giải gắn với 2.3.2(b)}

\tp{Nhận xét} \TODO{xếp loại (phù hợp / phù hợp một phần / không phù hợp); nếu sai lệch thì nêu nguyên nhân}

\subsection{Đánh giá dự báo}

\noidung{Mốc chuẩn và thước đo sai số}
\tp{Nội dung lý thuyết cần kiểm chứng} \TODO{một câu mệnh đề cần kiểm chứng; dẫn mục 2.3.3(a)}
\tp{Cách triển khai} Dữ liệu thực. \textit{Bộ áp dụng:} Nile, AirPassengers, DAX (log giá). \TODO{quy trình cụ thể, tham số riêng của mục này}
\tp{Mã lệnh Python} \TODO{đoạn mã lõi 10--20 dòng (chỉ ký tự ASCII và chữ Việt, không dùng ký hiệu Hy Lạp trong mã); mã đầy đủ ở notebook}
\tp{Kết quả tính toán}
\input{\tabdir/tab_b12_11_ap_benchmarks.tex}
\input{\tabdir/tab_b12_11_dp_multistep.tex}
\input{\tabdir/tab_b12_11_dp_one_step.tex}
\input{\tabdir/tab_b12_11_ni_benchmarks.tex}
\TODO{chọn bảng/hình giữ lại, bỏ bớt phần phụ; viết câu dẫn}
\tp{Phân tích và diễn giải}
\bodl{(i) Nile.} \TODO{diễn giải gắn với 2.3.3(a)}

\bodl{(ii) AirPassengers.} \TODO{diễn giải gắn với 2.3.3(a)}

\bodl{(iii) DAX (log giá).} \TODO{diễn giải gắn với 2.3.3(a)}

\tp{Nhận xét} \TODO{xếp loại (phù hợp / phù hợp một phần / không phù hợp); nếu sai lệch thì nêu nguyên nhân}

\noidung{Gốc cuốn chiếu}
\tp{Nội dung lý thuyết cần kiểm chứng} \TODO{một câu mệnh đề cần kiểm chứng; dẫn mục 2.3.3(b)}
\tp{Cách triển khai} Dữ liệu thực. \textit{Bộ áp dụng:} Nile, AirPassengers, DAX (log giá). \TODO{quy trình cụ thể, tham số riêng của mục này}
\tp{Mã lệnh Python} \TODO{đoạn mã lõi 10--20 dòng (chỉ ký tự ASCII và chữ Việt, không dùng ký hiệu Hy Lạp trong mã); mã đầy đủ ở notebook}
\tp{Kết quả tính toán}
\input{\tabdir/tab_b12_12_ap_one_step.tex}
\input{\tabdir/tab_b12_12_ap_rolling_by_horizon.tex}
\input{\tabdir/tab_b12_12_dp_rolling.tex}
\input{\tabdir/tab_b12_12_ni_rolling.tex}
\figx{fig_b12_12_ap_rmse_by_horizon}{AP: RMSE gốc cuốn chiếu theo chân trời, airline so với naive mùa vụ}
\TODO{chọn bảng/hình giữ lại, bỏ bớt phần phụ; viết câu dẫn}
\tp{Phân tích và diễn giải}
\bodl{(i) Nile.} \TODO{diễn giải gắn với 2.3.3(b)}

\bodl{(ii) AirPassengers.} \TODO{diễn giải gắn với 2.3.3(b)}

\bodl{(iii) DAX (log giá).} \TODO{diễn giải gắn với 2.3.3(b)}

\tp{Nhận xét} \TODO{xếp loại (phù hợp / phù hợp một phần / không phù hợp); nếu sai lệch thì nêu nguyên nhân}

\noidung{Hàm mất mát và kiểm định Diebold–Mariano}
\tp{Nội dung lý thuyết cần kiểm chứng} \TODO{một câu mệnh đề cần kiểm chứng; dẫn mục 2.3.3(c)}
\tp{Cách triển khai} Dữ liệu thực. \textit{Bộ áp dụng:} AirPassengers, Nile. \textit{Không áp dụng cho:} DAX (log giá) \TODO{lý do ngắn}. \TODO{quy trình cụ thể, tham số riêng của mục này}
\tp{Mã lệnh Python} \TODO{đoạn mã lõi 10--20 dòng (chỉ ký tự ASCII và chữ Việt, không dùng ký hiệu Hy Lạp trong mã); mã đầy đủ ở notebook}
\tp{Kết quả tính toán}
\input{\tabdir/tab_b12_13_diebold_mariano.tex}
\input{\tabdir/tab_b12_13_loss_functional.tex}
\input{\tabdir/tab_b12_13_ranking.tex}
\TODO{chọn bảng/hình giữ lại, bỏ bớt phần phụ; viết câu dẫn}
\tp{Phân tích và diễn giải}
\bodl{(i) AirPassengers.} \TODO{diễn giải gắn với 2.3.3(c)}

\bodl{(ii) Nile.} \TODO{diễn giải gắn với 2.3.3(c)}

\tp{Nhận xét} \TODO{xếp loại (phù hợp / phù hợp một phần / không phù hợp); nếu sai lệch thì nêu nguyên nhân}

\noidung{Quy trình báo cáo dự báo}
\tp{Nội dung lý thuyết cần kiểm chứng} \TODO{một câu mệnh đề cần kiểm chứng; dẫn mục 2.3.3(d)}
\tp{Cách triển khai} Dữ liệu thực. \textit{Bộ áp dụng:} Nile, AirPassengers, DAX (log giá). \TODO{quy trình cụ thể, tham số riêng của mục này}
\tp{Mã lệnh Python} \TODO{đoạn mã lõi 10--20 dòng (chỉ ký tự ASCII và chữ Việt, không dùng ký hiệu Hy Lạp trong mã); mã đầy đủ ở notebook}
\tp{Kết quả tính toán}
\input{\tabdir/tab_b12_14_checklist.tex}
\figx{fig_b12_14_residual_acf}{ACF của phần dư (trái) và của phần dư bình phương (phải) cho NI, AP, DP}
\TODO{chọn bảng/hình giữ lại, bỏ bớt phần phụ; viết câu dẫn}
\tp{Phân tích và diễn giải}
\bodl{(i) Nile.} \TODO{diễn giải gắn với 2.3.3(d)}

\bodl{(ii) AirPassengers.} \TODO{diễn giải gắn với 2.3.3(d)}

\bodl{(iii) DAX (log giá).} \TODO{diễn giải gắn với 2.3.3(d)}

\tp{Nhận xét} \TODO{xếp loại (phù hợp / phù hợp một phần / không phù hợp); nếu sai lệch thì nêu nguyên nhân}
